\section{Results}\label{sec:results}

All results are aggregated across 10 seeds with $N = 1{,}000$ trials
per seed per scenario.  Wilson 95\% confidence intervals are reported
where applicable.  Significance ($p < 0.001$) is determined by
non-overlapping Wilson CIs between biological and naive variants.

\subsection{Metabolism: Clear Win for Priority Gating}

Table~\ref{tab:metabolism} summarizes the metabolism benchmark results.

\begin{table}[h]
\caption{Metabolism benchmark: Biological (ATP\_Store + MTORScaler) vs Ablated (ATP\_Store only) vs Naive (flat counter). N = total operations or pressure events across 10 seeds.}
\label{tab:metabolism}
\centering
\small
\begin{tabular}{llccccc}
\toprule
Scenario & Metric & Bio & Abl & Naive & $\Delta$ & N \\
\midrule
bursty & critical served & 1.000 & 1.000 & 0.398 & \textbf{+0.602} & 24{,}139 \\
bursty & ops completed & 0.719 & 0.715 & 0.685 & \textbf{+0.034} & 2{,}000{,}000 \\
gradual depl. & critical served & 1.000 & 1.000 & 0.769 & \textbf{+0.231} & 43{,}372 \\
gradual depl. & ops completed & 0.979 & 0.996 & 0.950 & \textbf{+0.029} & 2{,}000{,}000 \\
constant & ops completed & 1.000 & 1.000 & 1.000 & 0.000 & 2{,}000{,}000 \\
sudden spike & ops completed & 1.000 & 1.000 & 1.000 & 0.000 & 2{,}000{,}000 \\
mixed priority & ops completed & 1.000 & 1.000 & 1.000 & 0.000 & 2{,}000{,}000 \\
\bottomrule
\end{tabular}
\end{table}

The headline result is \texttt{critical\_served\_under\_pressure}.  Under
bursty load, the biological system served 100\% of critical operations
($\text{priority} \geq 5$) during resource pressure, while the naive
flat counter served only 39.8\% ($\Delta = +0.602$, significant).  Under
gradual depletion, the gap is smaller but still significant: 100\% vs
76.9\% ($\Delta = +0.231$).

Against the naive baseline, the biological system actually \emph{improves}
total throughput: 71.9\% versus 68.5\% under bursty load ($\Delta = +0.034$),
because priority gating preserves budget for later operations that the
flat counter would miss entirely.  The throughput tradeoff appears when
comparing biological against ablated: under gradual depletion, biological
completed 97.9\% versus ablated's 99.6\%, because the MTORScaler's
anticipatory conservation gates some non-critical operations that
ATP\_Store alone would have served.  This is the intended tradeoff:
the mTOR scaler trades 1.7 percentage points of throughput against the
ablated control for guaranteed critical service via feature gating.

Under constant load, sudden spike, and mixed priority, all three
variants perform identically.  The biological complexity only matters
under specific pressure patterns where resource contention forces
choices about what to serve.

\subsection{Quorum Sensing: Unique Precision--Recall Operating Point}

Table~\ref{tab:quorum} summarizes key quorum sensing results.

\begin{table}[h]
\caption{Quorum sensing benchmark: selected configurations. Bio = QuorumSensingBio (signal accumulation + decay), Abl = IndependentActors (no coordination), Naive = MajorityVote.}
\label{tab:quorum}
\centering
\small
\begin{tabular}{llccccc}
\toprule
Config & Metric & Bio & Abl & Naive & $\Delta$ & N \\
\midrule
static, $n$=20, $c$=8 & TPR & 0.866 & 1.000 & 0.005 & \textbf{+0.861} & 300{,}000 \\
static, $n$=10, $c$=4 & TPR & 0.711 & 1.000 & 0.052 & \textbf{+0.659} & 300{,}000 \\
static, $n$=20, $c$=4 & TPR & 0.748 & 1.000 & 0.000 & \textbf{+0.748} & 300{,}000 \\
gradual, $n$=20, $c$=8 & TPR & 0.718 & 0.999 & 0.002 & \textbf{+0.716} & 216{,}959 \\
\midrule
all $c$=0 configs & FPR & 0.000 & 0.048--0.180 & 0.000 & 0.000 & 300K ea. \\
\bottomrule
\end{tabular}
\end{table}

The three coordination strategies occupy distinct positions on the
precision--recall curve:

\begin{itemize}[nosep]
  \item \textbf{Independent actors} (ablated): 100\% TPR but 5--18\% FPR.
        Catches everything but cries wolf on 1 in 6--20 clean time steps.
  \item \textbf{Majority vote} (naive): 0\% FPR but near-0\% TPR.
        Almost never detects threats because clean agents ($\mu = 0.15$)
        always outvote compromised agents at the $0.5$ threshold.
  \item \textbf{Signal accumulation} (biological): 0\% FPR with
        71--87\% TPR at 40\% compromise.  Zero false alarms with
        meaningful detection.
\end{itemize}

The biological model's advantage comes from two structural properties:
(1)~\emph{continuous signals} average out noise (unlike binary votes),
and (2)~\emph{temporal decay} ensures stale evidence ages out (unlike
a vote that persists forever).  The combination naturally filters noise
while remaining sensitive to sustained patterns.

Detection sensitivity scales with both the number of agents and the
compromise fraction, consistent with the biological signal-accumulation
model: more agents producing more signal yield higher accumulated
concentrations relative to the threshold.

\subsection{Epiplexity: Honest Negative}

Table~\ref{tab:epiplexity} summarizes the epiplexity benchmark results.

\begin{table}[h]
\caption{Epiplexity benchmark: Biological (EpiplexityMonitor, two-signal Bayesian) vs Naive (cosine similarity + timeout). All results with MockEmbeddingProvider.}
\label{tab:epiplexity}
\centering
\small
\begin{tabular}{llccccc}
\toprule
Scenario & Metric & Bio & Abl & Naive & $\Delta$ & N \\
\midrule
loop & accuracy & 0.467 & 0.080 & 0.940 & \textbf{$-$0.473} & 500{,}000 \\
loop & FP rate & 0.000 & 0.000 & 0.750 & \textbf{$-$0.750} & 40{,}000 \\
loop & FN rate & 0.580 & 1.000 & 0.000 & \textbf{+0.580} & 460{,}000 \\
\midrule
trophic w. & accuracy & 0.431 & 0.400 & 0.480 & $-$0.049 & 500{,}000 \\
trophic w. & FN rate & 0.948 & 1.000 & 0.866 & +0.081 & 300{,}000 \\
\midrule
false stagn. & accuracy & 0.540 & 1.000 & 0.993 & \textbf{$-$0.453} & 500{,}000 \\
false stagn. & FP rate & 0.052 & 0.000 & 0.007 & +0.045 & 500{,}000 \\
\midrule
convergence & accuracy & 0.473 & 1.000 & 1.000 & \textbf{$-$0.527} & 500{,}000 \\
exploration & accuracy & 1.000 & 1.000 & 1.000 & 0.000 & 500{,}000 \\
\bottomrule
\end{tabular}
\end{table}

The naive cosine-similarity detector outperforms the biological
Bayesian two-signal design on detection accuracy across all scenarios,
including the pathway-grounded trophic-withdrawal scenario that was
specifically designed to favor the biological approach.

However, the false-positive/false-negative breakdown reveals a more
nuanced picture.  On the loop scenario:

\begin{itemize}[nosep]
  \item \textbf{Biological}: 0\% FP rate, 58\% FN rate.
        Never false-alarms, but misses most stagnation.
  \item \textbf{Naive}: 75\% FP rate, 0\% FN rate.
        Catches all stagnation, but false-alarms on 3 of 4 healthy steps.
\end{itemize}

This is the same precision--recall tradeoff observed in quorum
sensing---but here the biological variant's recall (42\%) is too low
to be operationally useful, while its precision advantage (zero false
positives) is less valuable for stagnation detection than for threat
consensus.

The trophic-withdrawal scenario---designed to produce the biological
signature of epistemic atrophy (declining novelty with maintained
perplexity)---shows marginal differences: naive 48.0\% accuracy versus
biological 43.1\%.  With mock embeddings, the novelty signal does not
carry enough semantic information to make the two-signal distinction
meaningful (Section~\ref{sec:limitations}).
